
\minitoc 

% ===================================================================================================================
% Introduction 



% ===================================================================================================================
% Contexte

\section{Contexte}

Dans cette section, nous posons le cadre de l'étude des statistiques ainsi que l'ensemble 
des notations utilisées. 

\subsection{Différents Modèles}

L'étude des statistiques se résume à l'étude de la loi de variables aléatoires ou de suites de 
variables aléatoires. Nous devons donc définir un cadre dans lequel ces différents objets vivrons. 
On considère ainsi des variables aléatoires $X$ définies telles que : 
\begin{equation}
    X : 
        \begin{cases}
            ( \Omega, \mathcal{T}, \myP) \longrightarrow (\mathcal{X}, \mathcal{A}, \myP^X) \\ 
            \omega \longmapsto X( \omega)
        \end{cases}    
\end{equation}
Dans notre contexte, nous n'avons accès qu'à des \emph{réalisations} $X( \omega)$ de la variable 
aléatoire $X$ et nous souhaitons déterminer sa loi $ \myP^X$. 
Pour cela, nous faisons une hypothèse appelée \emph{hypothèse paramtérique forte} qui suppose que 
la loi de $X$ ne dépend que d'un paramètre $ \theta \in \Theta$, où $ \Theta \subset \R^d$. 
Ainsi, la recherche de la loi $ \myP^X$ se réduit à la détermination d'un paramètre $ \theta_*$ tel 
que $ \myP^X = \myP_{\theta_*}$. 

\begin{definition}[Modèle Paramétrique]
    On appellera \emph{modèle paramétrique} l'espace probabilisé $( \mathcal{X}, \mathcal{A}, \mathcal{P}_\Theta)$ 
    tel que $ \mathcal{P}_\Theta = \{ \myP_\theta \mid \theta \in \Theta\}$ où $ \Theta$ est un ouvert convexe 
    de $ \R^d$. 
\end{definition}

Or cet espace ne nous permet de considérer que de simples variables aléatoires. Dans un cadre plus général, 
nous souhaiterions étudier des \emph{échantillons} $(X_1, \dots, X_n)$ iid d'une variable aléatoire 
$X$. Nous définissons pour cela le \emph{modèle d'échantillonnage}. 

\begin{definition}[Modèle d'échantillonnage]
    Soit $( \mathcal{X}, \mathcal{A}, \mathcal{P}_{\Theta})$ un modèle paramétrique. 
    Un \emph{échantillon} $(X_1( \omega), \dots, X_n( \omega))$ de ce modèle est un $n$-uplet de réalisations 
    indépendantes d'une variable aléatoire $X$ de loi $ \myP_\theta \in \mathcal{P}_\Theta$. 
    On note $ \myP^\{X_1, \dots, X_n\} = \myP_{\theta_*}^{n}$ la loi de cet échantillon. Le modèle 
    paramétrique de cet échantillon est donc : 
        $$ ( \mathcal{X}, \mathcal{A}, \mathcal{P}_\Theta)^n = ( \mathcal{X}^n, \mathcal{A}^n, \mathcal{P}_{\Theta}^{n}) $$
\end{definition}

On a donc : 
    $$ 
        (X_1, \dots, X_n) : 
            \begin{cases}
                ( \Omega, \mathcal{T}, \myP) &\longrightarrow ( \mathcal{X}, \mathcal{A}, \mathcal{P}_\Theta)^n \\ 
                \omega &\longmapsto (X_1( \omega), \dots, X_n( \omega))
            \end{cases}
    $$

Enfin, pour l'étude plus théorique de tels objets, il serait intéressant de pouvoir faire 
tendre vers l'infini la taille de ces échantillons $(X_1, \dots, X_n)$. On définit donc un 
modèle dit \emph{asymptotique} : 
    $$ 
        (X_i)_{i \in \N^*} : 
        \begin{cases}
            ( \Omega, \mathcal{T}, \myP) & \longrightarrow ( \mathcal{X}, \mathcal{A}, \mathcal{P}_\Theta)^\infty \\ 
            \omega & \longmapsto \{X_i( \omega) \mid i \in \N_*\}
        \end{cases}
    $$

\subsection{Les variables aléatoires d'un modèle}

L'objet principal que nous allons manipuler permet d'effectuer des opérations sur des échantillons 
d'une variable aléatoire. On les appelle les \emph{statistiques}, notées $S$. 

\begin{definition}[Statistique]
    Une fonction mesurable $S$ est une \emph{statistique} si elle ne dépend pas 
    de $ \theta$. On a donc : 
        $$ 
            S(X_1, \dots, X_n) : 
                \begin{cases}
                    ( \mathcal{X}, \mathcal{A})^n \longrightarrow ( \mathcal{Y}, \mathcal{B}) \\ 
                    (X_1, \dots, X_n) \longmapsto S(X_1, \dots, X_n)
                \end{cases}
        $$ 
    On notera aussi $ \myP_\theta^S = ( \myP_\theta^n)^S$ la mesure image de 
    $\myP_\theta^n$ par $S$. 
\end{definition}

Cette translation de nos observations $(X_1, \dots, X_n)$ implique la définition d'un \emph{modèle image} 
nous permettant de manipuler ce nouvel objet $S(X_1, \dots, X_n)$. 

\begin{definition}[Modèle Image]
    Le modèle image de $( \mathcal{X}, \mathcal{A}, \mathcal{P}_\Theta)^n$ par $S_n$ 
    est donc $ ( \mathcal{Y}, \mathcal{B}, \mathcal{P}_\Theta^{S_n})$, avec : 
        \begin{equation}
            \mathcal{P}_\Theta^{S_n} := \{ \myP_\theta \circ S_n^{-1} \mid \theta \in \Theta\}
        \end{equation}
\end{definition}

Pour rappel, l'objectif des statistiques est de déterminer le paramètre $ \theta_* \in \Theta$. 
Pour cela, il faut que la donnée utilisée ($(X_1, \dots, X_n)$, $X$, ou $S(X_1, \dots, X_n)$) contienne 
le plus d'information possible sur $ \theta$. 
Il serait donc préférable que la translation $S$ ne réduise pas cette quantité. De telles statistiques seront 
donc dites \emph{exhaustives}. Elles seront caractérisées par la loi conditionnelle de l'échantillon. 

\begin{definition}[Loi conditionnelle d'une statistique]
    Soit $T$ une statistique. La loi conditionnelle de $(X_1, \dots, X_n)$ sachant 
    $T(X_1, \dots, X_n)$ est définie par : 
        \begin{equation}
            \myP_\theta^{(X_1, \dots, X_n) \mid T}(A) = \myP((X_1, \dots, X_n) \in A \mid T) 
            = \mathbb{E}_\theta( \mathbbm{1}_{(X_1, \dots, X_n) \in A} \mid T )
        \end{equation}
    pour tout $A \in \mathcal{A}^n$. 
\end{definition}

\begin{definition}[Statistique Exhaustives]
    Soit $T$ une statistique sur $( \mathcal{X}, \mathcal{A}, \mathcal{P}_\Theta)^n$ est dite \emph{exhaustive en $ \theta$}
    si $ \myP_\theta^{(X_1, \dots, X_n) \mid T}$ ne dépend pas de $ \theta$. Cela revient 
    à dire que l'application : 
        $$ \theta \longmapsto \frac{ \myP_\theta^n (A \cap T^{-1}(C))}{ \myP_\theta^T(C)} 
            = \frac{ \myP_\theta((X_1, \dots, X_n) \in A \cap T(X_1, \dots, X_n) \in C)}{ \myP_\theta T(X_1, \dots, X_n) \in C} $$
    est constante pour tout $ \theta \in \{ \theta \in \Theta \mid \myP_\theta^T(C) > 0\}$. 
\end{definition}

Une statistique exhaustive contient toute l'information nécessaire sur $ \theta^*$ contenue 
dans l'échantillon. C'est donc une bonne statistique à utiliser. 

\begin{proposition}
    $T$ est une \emph{statistique exhaustive} ssi pour toute statis $S$ réelle intégrable, on a : 
        \begin{equation}
            \mathbb{E}_\theta(S \mid \sigma(T)) = \varphi( T)
        \end{equation}
    (i.e si cela ne dépend pas de $ \theta$.)
\end{proposition}

\begin{example}[Statistique exhaustive]
    Soit $(X_1, \dots, X_n)$ un échantillon d'une variable aléatoire $ X \sim \mathcal{B}(1, \theta)$. Posons 
    $S_n = \sum_{i=1}^{n} X_i \sim \mathcal{B}(n, \theta)$. On a ici, $ \mathcal{X} = \{0, 1\}^n$. 
    On cherche à estimer $ \theta$. Montrons que $S$ est une statistique exhaustive. 
    Soit $k \in \llbracket 1, n \rrbracket$, on : 
    \begin{align*}
        & \myP_\theta ((X_1, \dots, X_n) = (x_1, \dots, x_n) \mid S_n = k) \\ 
        &= \frac{ \myP_\theta( (X_1, \dots, X_n) = (x_1, \dots, x_n) \cap S_n = k)}{ \myP_\theta(S_n = k)} \\ 
        &= \mathbbm{1}_{\left\{\sum_{i=1}^{n} x_i = k\right\}} \times  \frac{ \myP_\theta( (X_1, \dots, X_n) = (x_1, \dots, x_n))}{ \myP_\theta(S_n = k)} \\ 
        &= \mathbbm{1}_{\left\{\sum_{i=1}^{n} x_i = k\right\}} \times \prod_{i=1}^{n} \frac{ \myP_\theta(X_i = x_i)}{ \binom{n}{k} \theta^k (1 - \theta)^{n - k}} \\ 
        &= \frac{1}{\binom{n}{k}}
    \end{align*}
    qui est indépendant de $ \theta$ donc $S_n$ est bien une statistique exhaustive en $ \theta$. 
\end{example}

\begin{definition}[Statistique libre et Variable pivotable]
    On appelle \emph{statistique libre} une statistique $S$ telle que $ ( \myP_\theta^n)^S$ ne dépend 
    pas de $ \theta$. 
    On appelle \emph{variable aléatoire pivotable} $U_\theta(X_1, \dots, X_n)$ une fonction de $ \theta$
    et de l'échantillon telle que, sous la loi $ \myP_\theta$ pour $(X_1, \dots, X_n)$, la mesure 
    image $ ( \myP_\theta^n)^{U_\theta}$ ne dépend pas de $ \theta$. 
\end{definition}


% ===================================================================================================================
% Information Qualitative

\section{Information Qualitative}

\subsection{Domination et Densités}

\begin{definition}[Domination]
    On dira qu'une mesure $ \mu$ positive \emph{domine} $ \myP_\theta$ sur $( \mathcal{X}, \mathcal{A})$ si : 
    \begin{equation}
        \boxed{
            \forall A \in \mathcal{A}, \quad \mu(A) = 0 \Longrightarrow \myP_\theta(A) = 0 
        }
    \end{equation}
    On notera alors $ \myP_\theta \ll \mu$.
    Si c'est le cas pour tout $ \theta \in \Theta$, on dira alors que \emph{le modèle est dominé par $ \mu$}.
\end{definition}

\begin{lemma}
    Plus formellement, $ \myP_\theta \ll \mu$ ssi 
    \begin{equation}
        \forall \varepsilon > 0, \exists \delta > 0 \quad \text{tel que} \quad 
        \forall A \in \mathcal{A}, \mu(A) < \delta \Longrightarrow \myP_\theta(A) < \varepsilon 
    \end{equation}
\end{lemma}

\begin{theorem}[Radon-Nikodym]
    Soit $ \mu$ $ \sigma$-finie. On a $ \myP_\theta \ll \mu$ ssi 
    \begin{equation}
        \forall \theta \in \Theta, \exists f \in L^1_+( \mu) \quad \text{telle que} \quad 
        \forall A \in \mathcal{A}, \quad \myP_\theta(A) = \int_A f \; d \mu 
    \end{equation}
    La densite $f$ de $ \myP_\theta$ est unique à égalité $ \mu$-presque partout. Elle est notée 
    $f = d \myP_\theta / d \mu$ comme dérivée de Radon-Nicodym. 
\end{theorem}

\begin{remark}[Cas discret]
    Le cas des lois discrètes peut paraître contre-intuitif, mais s'explique bien 
    grâce à ce théorème. Soit $ \mathcal{X}$ fini ou dénombrable et $ \mathcal{A} = \mathcal{P}( \mathcal{X})$. 
    Soit $ m_c$ la mesure de comptage sur $ \mathcal{A}$. Par définition de l'intégrale de Lebesgue, pour toute 
    fonction $ f : \mathcal{X} \longrightarrow \R$ et $ A \in \mathcal{A}$, on a : 
        $$ \int_A f d \; m_c = \sum_{x \in A} f(x) m_c(x) = \sum_{x \in A} f(x) $$ 
    Soit $ \mathcal{P}_\Theta$ une famille de lois sur $ \mathcal{X}$ et $n = 1$. On a $ \mathcal{P}_\theta \ll m_c$ et 
    les densités $f = d \myP_\theta / d m_c$ existent et vérifient :
        $$ \forall A = {x} \in \mathcal{A}, \quad \myP_\theta(\{x\}) = \int_{\{x\}} f \; d m_c = f(x) $$
    On a donc bien $f(x) = \myP_\theta(X = x)$. 
\end{remark}

\subsection{Hypothèses de Cramer-Rao}

Pour tous les raisonnements suivants sur notre modèle, nous devons définir des \emph{hypothèses}. 
Celles-ci nous assureront certaines propriétés nécessaires à l'application des théorèmes suivants. 
Elles sont donc : 
    \begin{center}
        $(H_0)$ : $ \mathcal{P}_\Theta$ est identifiable et dominé. \\
        $(H_1)$ : $ \mathcal{P}_\Theta$ est identifiable et homogène. 
    \end{center}

\begin{proposition}
    S'il existe $ \mu$ $ \sigma$-finie telle que $ \mathcal{P}_\Theta \ll \mu$ et $ \mu(\{f_{\theta_0} \not = f_{\theta_1}\}) > 0$ 
    pour tout $ \theta_0 \not = \theta_1$, alors le modèle $( \mathcal{X}, \mathcal{A}, \mathcal{P}_\Theta)$ 
    est dit \emph{identifiable}. 
\end{proposition}

\begin{definition}[Modèle homgène]
    Le modèle $( \mathcal{X}, \mathcal{A}, \mathcal{P}_\Theta)$ est dit \emph{homogène} si 
    $ \mathcal{P}_\Theta \ll \mu$ pour tout $ \theta \in \Theta$. 
\end{definition}

\begin{proposition}
    Le modèle $( \mathcal{X}, \mathcal{A}, \mathcal{P}_\Theta)$ est homogène ssi, il existe $ \mu$ 
    $ \sigma$-finie telle que : 
    \begin{enumerate}[label=\roman*)]
        \item $ \mathcal{P}_\Theta \ll \mu$ 
        \item $\{f_\theta > 0\}$ $ \mu$-presque partout pour tout $ \theta \in \Theta$. 
    \end{enumerate}
\end{proposition}

En conséquence, pour tout $ \theta \in \Theta$, $ \mathcal{P}_\Theta \ll \myP_\theta$ donc 
les $ \myP_\theta$ ont tous les mêmes négligeables et la densité de $ \myP_{ \theta_2}$ par rapport 
à $ \myP_{ \theta_1}$ s'écrit : 
    $$ 
        \frac{ \partial \myP_{ \theta_2}}{ \partial \myP_{ \theta_1}} = \frac{ f_{ \theta_2}}{f_{ \theta_1}} 
        = \frac{ \partial \myP_{ \theta_2} / d \mu}{ \partial \myP_{ \theta_1} / d \mu}
    $$ 
En effet, on a pour tout $A \in \mathcal{A}$ : 
    \begin{align*}
        \myP_{ \theta_2}(A) &= \int_A \frac{f_{ \theta_2}}{f_{ \theta_1}} \; d \myP_{ \theta_1} 
        = \int_A \frac{f_{ \theta_2}}{f_{ \theta_1}} f_{ \theta_1} \; d \mu \\ 
        &= \int_A f_{ \theta_2} \; d \mu = \myP_{\theta_2}(A). 
    \end{align*}

\subsection{Vraissemblance}

Venons à la définition principale de ce cours : la fonction de vraissemblance, notée $ L_\theta$ 
pour "likelihood". Elle correspond à la densité de l'échantillon $(X_1, \dots, X_n)$ a posteriori aux 
résultats $X_i( \omega)$. 

\begin{definition}[Fonction de vraissemblance]
    Dans un modèle vérifiant $H_0$ et $H_1$, on appelle \emph{fonction de vraissemblance}
    l'application $L$ définie par : 
        $$ 
            L : 
                \begin{cases}
                    \Theta \times \mathcal{X}^n \longrightarrow (\R_+, \mathcal{B}_{\R_+}) \\ 
                    \theta, X_1, \dots, X_n \longmapsto L( \theta, X_1, \dots, X_n)
                \end{cases}
        $$ 
    elle vérifie : 
        $$
            L(\theta, X_1, \dots, X_n) = \prod_{i=1}^{n} f_\theta(X_i) 
        $$
\end{definition}

Cette nouvelle application nous permet de caractériser plus simplement la notion de statistique 
exhaustive. 

\begin{theorem}[Critère de Fisher-Neyman]
    $T$ est une statistique exhaustive ssi 
        $$ \forall \theta \in \Theta, \quad L( \theta, X_1, \dots, X_n) = h(X_1, \dots, X_n) \times g( \theta, T(X_1, \dots, X_n)) $$
    avec $ h : \mathcal{X}^n \longrightarrow \R$ et $ g : T( \mathcal{X}^n) \longrightarrow \R$ positives mesurables. 
\end{theorem}

\subsection{Minimalité}

On souhaite définir des statistiques dites minimales, contenant uniquement l'information nécessaire 
sur $ \theta$ à partir de l'échantillon. Pour cela, on pose : 
    $$ \mathcal{A}^* := \bigcap_{T \text{ exhaustives}} \sigma(T) $$ 
appelée la \emph{tribu exhaustive minimale}. Elle contient donc toute l'information apportée par 
$(X_1, \dots, X_n)$ sur $ \theta$. 

\begin{definition}[Statistique exhaustive minimale]
    On dira d'une statistique exhaustive $S^*$ qu'elle est \emph{minimale} si $ \sigma(S^*) = \mathcal{A}^*$. 
\end{definition}

\begin{proposition}[Caractérisation de $ \mathcal{A}$]
    Sous $ H_1$, pour $ \theta_0 \in \Theta$ fixé, on a : 
        $$ 
            \mathcal{A} = \sigma \left\{
                \frac{L_\theta}{L_{\theta_0}} (X_1, \dots, X_n) \mid \theta \in \Theta
            \right\}
        $$
\end{proposition}

\begin{proposition}[Statistiques exhaustives minimales]
    Soit $S$ une statistique exhaustive telle que : 
        \begin{equation}
            S(x_1, \dots, x_n) = S(y_1, \dots, y_n) \quad \iff \frac{L_\theta(x_1, \dots, x_n)}{L_\theta(y_1, \dots, y_n)}
            \text{ indépendants de } \theta
        \end{equation}
    alors $S$ est \emph{minimale}. 
\end{proposition}

\begin{example}[Loi exponentielle]
    Soit la densité de la loi exponentielle $ f_\theta(x) = \theta e^{ - \theta x} $. On cherche à déterminer $S^*$. 
    Posons $ S_n = \sum_{i=1}^{n} X_i$. On a : 
        $$ L_\theta(X_1, \dots, X_n) = \prod_{i=1}^{n} f_\theta(X_i) = \theta^n e^{- \theta S_n} $$ 
    Pour tout $ \theta_0 \not = \theta_1$, on a :
        $$ 
            \frac{L_{\theta_1}}{L_{\theta_2}}(X_1, \dots, X_n) = \exp (( \theta_0 - \theta_1) S_n + n(\log \theta_1 - \log \theta_0))
        $$ 
    Donc $S$ est fonction mesurable du rapport $ L_{\theta_1}/L_{\theta_2}$, elle est donc minimale. 
\end{example}


% ===================================================================================================================
% Information de Küllback-Leibler

\section{Information de Küllback-Leibler}

\begin{definition}[Contraste de Küllback-Leibler]
    Le \emph{contraste de Küllback} est défini par : 
    \begin{equation}
        K( \theta_0 \mid \theta_1) = \mathbb{E}_{\theta_0} \left(\log \frac{f_{\theta_0}}{f_{\theta_1}} \right) = 
        \int_\mathcal{X} \log \left( \frac{f_{\theta_0}}{f_{\theta_1}}\right) f_{\theta_0} \; d \mu 
    \end{equation}
\end{definition}

On parle de \emph{contraste}, car cette application n'est pas une distance par manque de symmétrie. 

\begin{prop}[Modèle dominé]
    Dans tout modèle dominé $ \mathcal{P}_\Theta << \mu$, $K( \theta_1 \mid \theta_2)$ est bien définie. 
\end{prop}

\begin{proposition}
    L'information de Küllback $ K( \theta_0 \mid \theta_1)$ ne dépend pas de la mesure dominante $ \mu$. 
    On a aussi $ K( \theta_0 \mid \theta_1) \geqslant 0$, et $ K( \theta_0 \mid \theta_1) = 0$ ssi 
    $ \myP_{\theta_0} = \myP_{\theta_1}$. 
\end{proposition}






